\section{Introduction}

\begin{figure*}[t]
\centering
\includegraphics[width=\textwidth]{figures/figure1.pdf}
\caption{Why escalation needs a guarantee. (a) The team costs up to 6.3 times a single reasoning agent.
(b) Some tasks the team solves are already solved by the direct answer $L_0$ (fill) or by $L_0$ or a
single agent $L_1$ (outline). (c) Confidence only partly separates right from wrong, and on three settings ranks them backwards. (d) A fixed threshold
that spends as much on average as \climb{} exceeds the tolerance on some settings, while \climb{} stays within a tolerance of $\eps = 5\%$ (replay).}
\label{fig:motivation}
\end{figure*}

Multi-agent systems built from large language models divide a task among roles: a planner decomposes it,
an executor solves it and a reviewer critiques the result over several rounds
\citep{du2024debate,wu2024autogen,hong2024metagpt}. Such pipelines are costly and not uniformly better. On
the three open models we study, a plan--execute--review pipeline spends 5.4k to 175k tokens per task, 1.3
to 6.3 times the cost of a single agent that reasons about the task once
(Figure~\ref{fig:motivation}a), yet some of the tasks it solves are already solved by a direct answer or a
single agent (Figure~\ref{fig:motivation}b). The team can also reduce accuracy: on some tasks a direct answer
or a single agent is correct and the team is not, mostly because its reviewer fails to correct a wrong first
draft, and almost half of these losses against the direct answer recur when the team is rerun
(\S\ref{sec:analysis}). The relevant question is thus less whether collaboration helps on average than
\emph{when it is worth its cost} for a given input.

Invoking the team only when a cheaper configuration is likely to fail can reduce cost, but the difficulty lies in
the rule that decides when. At inference time the system observes only weak proxies for the quality of a cheap
answer: the mean token log-probability of a direct answer separates correct from incorrect answers with an
AUROC of 0.22 to 0.81, falling below chance where greedy decoding produces confident repetitions
(Figure~\ref{fig:motivation}c). Such proxies can work on average but offer no guarantee. A fixed confidence
threshold with the same average cost as our method loses more than 2 points of accuracy on seven of our twelve
settings and up to 16 (Figure~\ref{fig:motivation}d), and always stopping at the single agent loses up to 9.
Without a bound on the accuracy given up, such savings are difficult to trust, and practitioners may
reasonably continue to pay for the full pipeline.

We instead minimise expected cost subject to a user-set bound $\eps$ on the accuracy lost relative to always running
the pipeline. A task climbs a ladder whose rungs are cheaper configurations of the same model, a direct answer
$L_0$ and then a single agent $L_1$ at the pipeline's own reasoning setting. It stops at a rung once the
evidence suffices and escalates to the full pipeline otherwise. Because skipping the team can also retain
correct answers the team would miss, this loss is \emph{signed}. When the gain from escalation decreases with
the confidence of the cheap answer, the optimal rule is a confidence threshold, and the safe thresholds are the
most conservative ones. \climb{} exploits this ordering without assuming the condition: following
Learn-then-Test \citep{angelopoulos2021ltt}, it tests the thresholds from conservative to aggressive with a
betting test on the signed loss \citep{waudbysmith2024betting}. This yields a distribution-free guarantee: with
probability at least $1-\delta$ for a chosen error level $\delta$, the deployed policy loses at most $\eps$ accuracy relative to the team.
\climb{} answers cheaply where the guarantee permits and runs the pipeline elsewhere, including on every task
when no cheaper policy can be certified, so a weak or misleading signal tends to reduce savings rather than
accuracy.

Across three open models and four benchmarks, \climb{} stayed within its tolerance in over 96\% of i.i.d.\
calibration draws, where the guarantee requires 90\%, while saving 38\% of the pipeline's tokens on average at
a 5-point tolerance and 65\% at 10 points. Where its signal was misleading, it saved little rather than losing
accuracy. The choice of test can matter as much as the signal: on the same records, the Hoeffding--Bentkus
bound of \citet{wynn2026corrupted} certifies nothing cheaper than the pipeline, and clipping the loss at zero
certifies less on most settings. Our contributions are as follows:
\begin{itemize}
    \setlength{\itemsep}{1pt}
    \setlength{\parskip}{0pt}
    \setlength{\parsep}{0pt}
    \item We measure what a plan--execute--review team adds over a single agent, including the correct
    answers it loses and where those losses arise, and cast the decision of when to call it as a
    constrained problem with a signed loss, under which one monotone-gain condition makes a threshold optimal and testing in order lossless (\S\ref{sec:motivation}, \S\ref{sec:structure}).
    \item We introduce \climb{}, a training-free procedure that escalates on confidence and on the
    agreement of two cheap rungs, and certifies the accuracy lost relative to the team with a
    variance-adaptive betting test inside fixed-sequence testing (\S\ref{sec:calib}).
    \item On three models and four benchmarks we verify that the guarantee holds in practice, show that uncertified rules of
    similar cost can lose unpredictable amounts, and find that, when most losses are zero, variance adaptivity matters as much as keeping the loss's sign (\S\ref{sec:results}).
\end{itemize}
